\section{Die Liouvilleschen Sätze}

\begin{bla}{Erster Liouvillscher Satz}\label{2.1}
Sei $f : \mathbb{C} \to \mathbb{C}$ eine ganze Funktion.
Wenn $f$ elliptisch ist, dann ist $f$ konstant.
\end{bla}

\begin{proof}
Sei $f$ eine elliptische Funktion mit den Perioden $\omega_1$ und $\omega_2$.
Es genügt $f$ auf einem seiner Periodenparallelogramme zu betrachten.
Also betrachtet man die kompakte Menge
\begin{align*}
P := \lbrace t \cdot \omega_1 + r \cdot \omega_2 \ : \ t,r \in [0,1] \rbrace.
\end{align*}
und wählt $z \in P$ beliebig. Da $f$ holomorph auf $\mathbb{C}$ ist, ist $f$ stetig auf $P$.
Mit
\begin{align*}
| f(z) | \leq \max_{z \in P} |f(z)| < \infty.
\end{align*}
ist $f$ auf $P$ beschränkt. Aus der Periodizität folgt die Beschränkheit auf ganz $\mathbb{C}$.
Nach dem Satz von Liouville aus Analysis 3 ist $f$ damit konstant.
\end{proof}



\begin{bla}{Folgerungen}\label{2.2}
Aus dem vorherigen Satz ergeben sich die Folgerungen:
\begin{enumerate}%[leftmargin=*]
\item[\textbf{(1)} ]
Jede nicht konstante elliptische Funktion enthält in jedem Periodenparallelogramm mindestens 
einen Pol.
\item[\textbf{(2)}]
Seien $f $ und $g$ elliptische Funktionen mit den gleichen Perioden.
Zusätzlichen seien deren Polstellen und zugehörige Hauptteile identisch.
Dann ist $f-g $ konstant.
\item[\textbf{(3)}]
Seien $f $ und $g$ elliptische Funktionen mit den gleichen Perioden.
Nun seien die Null-und Polstellen und deren Ordnungen gleich.
Dann ist $\nicefrac{f}{g} $ konstant.
\end{enumerate}
\end{bla}

\begin{proof}
\begin{enumerate}%[leftmargin=*]
\item[(1)]
Mit der Annahme, dass $f$ keine Polstellen besitzt, erhält man mit \ref{2.1}
einen Widerspruch. 
\item[(2)]
Seien $f,g$ elliptische Funktionen mit den angegebenen Voraussetzungen. Dann ist auch $f-g$ elliptisch.
Nun sei $z_0$ eine beliebige Polstelle von $f$ und $g$ mit Ordnung $K$. Dann lassen sich $f,g$ in $z_0$ als Laurentreihe darstellen. Es gilt also:
\begin{align*}
f(z)-g(z)
= &\sum_{k=1}^K a_{-k} \cdot (z-z_0)^{-k} + \sum_{k=0}^\infty a_k \cdot (z-z_0)^k
- \sum_{k=1}^K b_{-k} \cdot (z-z_0)^{-k} - \sum_{k=0}^\infty b_k \cdot (z-z_0)^k\\
= &\sum_{k=0}^\infty (a_k -b_k) \cdot (z-z_0)^k
\end{align*} 
%Da die Haupteile gleich sind gilt: $\forall k \in \lbrace 1,...,K \rbrace : a_{-k} = b_{-k}$
Damit besitzt $f-g$ an $z_0$ eine hebbare Singularität und da $z_0$ beliebig gewählt wurde ist $f-g$ auf ganz $\mathbb{C}$ holomorph fortsetzbar. %\\
Also ist $f-g$ mit \ref{2.1} konstant.  
\item[(3)]
Seien $f,g$ elliptische Funktionen mit den angegebenen Voraussetzungen. Dann ist auch $\nicefrac{f}{g}$ elliptisch.
Nun kürzen sich Pole und Nullstellen gegenseitig weg. Damit besitzt $\nicefrac{f}{g}$
keine Pole.
Somit ist $\nicefrac{f}{g}$ mit \ref{2.1} konstant.
\end{enumerate}
\end{proof}

\begin{bla}{Zweiter Liouvillscher Satz}\label{2.3}
Sei $f$ eine elliptische Funktion mit den Perioden $\omega_1,\omega_2$.
Weiter sei $S$ die Menge der Polstellen innerhalb eines Periodenparallelogramms von $f$. Dann gilt
\begin{align*}
\sum \limits_{z \in S } \Res(f,z) = 0.
\end{align*} 
\end{bla}

\begin{proof}
Zunächst platziert man ein Periodenparallelogramm,
so dass keine Polstellen von $f$ auf dessen Rand liegen.
Nun betrachtet man
\begin{align*}
\gamma_1 &:= [z_0,z_0+\omega_1]\\
\gamma_2 &:= [z_0+\omega_1,z_0+\omega_1+\omega_2]\\
\gamma_3 &:= [z_0+\omega_1+\omega_2,z_0+\omega_2]\\
\gamma_4 &:= [z_0+\omega_2,z_0]
\end{align*}
und erhält den Weg 
$\gamma := \gamma_1 \cup \gamma_2 \cup \gamma_3 \cup \gamma_4$,
welcher das verschobene Parallelogramm berandet.
Zusätzlich fällt mit der Periodizität von $f$ auf, dass 
\begin{align*}
\int \limits_{\gamma_3} f(z) \ \td{z} = -\int \limits_{\gamma_1} f(z) \ \td{z}
\ \text{und} \
\int \limits_{\gamma_2} f(z) \ \td{z} = -\int \limits_{\gamma_4} f(z) \ \td{z}
\end{align*}
gilt.
Also folgt mit dem Residuensatz, dass
\begin{align*}
\int \limits_\gamma f(z) \ \td{z}
&= 2 \pi  \mathrm{i}  \cdot \sum \limits_{z \in S } \Res(f,z)\\
&= \int \limits_{\gamma_1} f(z) \ \td{z}
+ \int \limits_{\gamma_2} f(z) \ \td{z}
+ \int \limits_{\gamma_3} f(z) \ \td{z}
+ \int \limits_{\gamma_4} f(z) \ \td{z}\\
&= \int \limits_{\gamma_1} f(z) \ \td{z}
- \int \limits_{\gamma_1} f(z) \ \td{z}
- \int \limits_{\gamma_4} f(z) \ \td{z}
+ \int \limits_{\gamma_4} f(z) \ \td{z}
= 0
\end{align*}
ist. 
Somit gilt die Behauptung.
\end{proof}

\begin{bla}{Folgerung}\label{2.4}
Es existiert keine elliptische Funktion $f$ mit $\ord(f) = 1$.
\end{bla}
\begin{proof}
Angenommen eine solche Funktion $f$ würde existieren.
Dann gibt es innerhalb eines Periodenparallelogramms
genau einen Pol $z_0$ erster Ordnung. Also gilt $\Res(f,z_0) \neq 0$.
Dies steht im direktem Widerspruch zum letzten Satz.
\end{proof}

\begin{bla}{Dritter Liouvillscher Satz}\label{2.5}
Sei $f$ eine elliptische Funktion mit $\ord(f) = n \geq 2$. In einem beliebig verschobenen Periodenparallelogramm $P_{z_0}$ gilt für ein festes $a \in \mathbb{C}$
\begin{align*}
\# \lbrace z \in P_{z_0} \ | \ f(z) - a = 0   \rbrace  = n,
\end{align*}
wobei bezüglich der Vielfachheit gezählt wird.
Das heißt, in $P_{z_0}$ wird jeder Funktionswert $a \in \mathbb{C}$ genau $n$-mal 
angenommen.
\end{bla}

\begin{proof}
Sei $a \in \mathbb{C}$ beliebig aber fest.
Nun betrachtet man $f(z) - a$ und verschiebt ein Periodenparallelogramm so, 
dass keine Urbilder von $a$ auf dessen Rand liegen.
Dazu wählt man dann den Weg $\gamma$ aus \ref{2.4} .
Mit \ref{1.4} ist auch $\nicefrac{f^\prime(z)}{f(z)-a}$ eine elliptische Funktion.
Es folgt mit dem Null- und Polstellen zählenden Integral,
\begin{align*}
\frac{1}{2 \pi \mathrm{i}} \cdot \int \limits_\gamma \frac{f^\prime(z)}{f(z)-a} \ \td{z}
= N - P \overset{\ref{2.4}}{=} 0
\quad \Longrightarrow \quad N = P
\end{align*}
wobei $N$  die Anzahl der Nullstellen von $f(z) - a$, also die Anzahl der Urbilder von $a$ unter $f$ ist und $P = n$
gilt.
Damit wird $a$ genau $n$-mal innerhalb eines Periodenparallelogramms angenommen.
\end{proof}

\begin{genericdef}{Bemerkung}\label{2.6}
Seien die Voraussetzungen wie im letzten Satz.
Dann entspricht die Anzahl der Nullstellen von $f$ (mit Vielfachheiten gezählt) der Ordnung von $f$.
Also gilt $ \ord f = \ord \nicefrac{1}{f}$.
\end{genericdef}